\begin{abstract}
Multi-agent LLM systems are often overseen by monitors that read their transcripts. We ask whether agents need a shared secret key to communicate in a way such monitors can neither read nor detect, and if not, what replaces the key. We split ``shared secret'' into separate ingredients: a key, common knowledge of the protocol, public randomness, channel entropy, shared weights, and each agent's private coins. We prove which ingredients each goal needs against monitors of different strength. A key is not necessary. However, a simple simulation argument shows that any convention the agents reach tacitly from shared knowledge gives them \emph{zero} advantage over a monitor with the same model and knowledge. Secrecy instead requires interaction, private randomness, entropy, prior common knowledge of the protocol, and a passive monitor. We test these predictions with open-weight models. With no agreed scheme, a \qwenbig partner recovers a hidden phrase 48\% of the time and a same-model auditor 50\% (McNemar $p=0.75$). A toy keyless covert key exchange run through the sampler of \qwenchan agrees on a 254-bit key in 10/10 sessions. An auditor that knows everything except the agents' private coins decodes 0/10 payloads and cannot detect the channel (NLL AUROC 0.46, 95\% CI [0.40, 0.53]). The same channel is fragile: an active warden that resamples only 0.25\% of tokens cuts key agreement to 2/10, because it desynchronizes the entropy-gated bit stream. Monitors should therefore share the agents' weights and context, and they should act on the channel instead of only reading it.
\end{abstract}
